% D52  v6 Van der Pol
% Auto-generated from markdown source
% origin: ai

\documentclass[12pt, a4paper]{article}
\usepackage[utf8]{inputenc}
\usepackage[T1]{fontenc}
\usepackage{amsmath, amssymb, amsfonts, amsthm}
\usepackage{textcomp}
\usepackage{geometry}
\usepackage{hyperref}
\usepackage{booktabs}
\usepackage{graphicx}
\usepackage{fancyhdr}
\usepackage{tcolorbox}
\usepackage{enumitem}

\geometry{margin=1in}
\pagestyle{fancy}
\fancyhf{}
\rhead{D52  v6 Van der Pol}
\lhead{\thepage}

\newtheorem{definition}{Definition}[section]
\newtheorem{theorem}{Theorem}[section]
\newtheorem{corollary}{Corollary}[section]
\newtheorem{proposition}{Proposition}[section]
\newtheorem{lemma}{Lemma}[section]
\newtheorem{conjecture}{Conjecture}[section]

\newcommand{\calI}{\mathcal{I}}
\newcommand{\calP}{\mathcal{P}}
\newcommand{\calE}{\mathcal{E}}
\newcommand{\calS}{\mathcal{S}}
\newcommand{\calH}{\mathcal{H}}
\newcommand{\calO}{\mathcal{O}}
\newcommand{\calF}{\mathcal{F}}
\newcommand{\calU}{\mathcal{U}}
\newcommand{\calV}{\mathcal{V}}
\newcommand{\calA}{\mathcal{A}}
\newcommand{\calG}{\mathcal{G}}
\newcommand{\calB}{\mathcal{B}}
\newcommand{\calC}{\mathcal{C}}
\newcommand{\calL}{\mathcal{L}}
\newcommand{\calM}{\mathcal{M}}
\newcommand{\calX}{\mathcal{X}}
\newcommand{\calY}{\mathcal{Y}}
\newcommand{\calW}{\mathcal{W}}
\newcommand{\calT}{\mathcal{T}}
\newcommand{\calZ}{\mathcal{Z}}
\newcommand{\calR}{\mathcal{R}}
\newcommand{\calQ}{\mathcal{Q}}
\newcommand{\calD}{\mathcal{D}}
\newcommand{\calN}{\mathcal{N}}
\newcommand{\calK}{\mathcal{K}}

\newcommand{\Tr}{\text{Tr}}
\newcommand{\E}{\mathbb{E}}
\newcommand{\diff}{\mathrm{d}}
\newcommand{\KL}{D_{\mathrm{KL}}}
\newcommand{\argmax}{\arg\max}
\newcommand{\argmin}{\arg\min}
\newcommand{\hbarF}{\hbar}

\newtcolorbox{thmbox}[1]{
  colback=blue!5!white, colframe=blue!75!black, fonttitle=\bfseries,
  title={#1}, breakable
}
\newtcolorbox{defbox}[1]{
  colback=green!5!white, colframe=green!75!black, fonttitle=\bfseries,
  title={#1}, breakable
}

\begin{document}

\title{D52  v6 Van der Pol}
\author{RALPH Loop}
\date{}
\maketitle

\begin{abstract}
This paper presents D52  v6 Van der Pol, formalizing the underlying principles using information-theoretic methods.

\textbf{Keywords}: D52  v6 Van der Pol, information theory, formalization.
\end{abstract}

\tableofcontents
\section{D52  v6 Van der Pol}


\subsection{D52 v3}


D52 ****
1. D52-D = 0.000000
2.  <5\%
3.  alpha

\subsection{v4v5v6}


\textbf{Van der Pol }

\[x'' - \mu(1-x^2)x' + x = \alpha\cos(\omega t)\]

\(\mu\)
\begin{itemize}
\item \(\mu = 0\)
\item \(\mu > 0\)
\end{itemize}


\#\#

\subsubsection{D52-A: - (v5)}


| \(\mu\) |  |  |  |
|-------|------|------|------|
| 0.0 | 0.50 | 1.00 | 0.1600 |
| 1.0 | 2.05 | 2.01 | 0.1520 |
| 5.0 | 2.47 | 2.02 | 0.0880 |
| 20.0 | 2.73 | 2.01 | 0.0280 |

\begin{itemize}
\item ~\(\mu\)  = -0.4284
\item (\(\mu \geq 5\))  = \textbf{-0.8587} -1
\end{itemize}


\subsubsection{D52-B:  alpha (v5)}


\begin{itemize}
\item alpha=0.02.05, alpha=1.51.00, alpha=5.05.73
\item alpha  1.5
\item alpha > 2 ~alpha
\end{itemize}


\subsubsection{D52-C:  (v5)}


\begin{itemize}
\item \textbf{Van der Pol  = 318.6} 0.78
\item \textbf{ = 0.0000}
\item alpha178324alpha=2.0
\end{itemize}


\subsubsection{D52-D: alpha-mu  (v5)}


\begin{itemize}
\item mu=0, alpha=0.7E=0.17
\item alpha  mu
\item pre-planning
\end{itemize}


\subsubsection{D52-F: - (v6)}


| \(\mu\) |  |  |  | rise/desc |  |
|-------|------|------|------|-----------|----------|
| 0.0 | 0.50 | 1.00 | 0.1600 | 0.5079 | 249.90 |
| 0.1 | 1.95 | 2.00 | 0.1600 | 0.4921 | 972.31 |
| 1.0 | 2.05 | 2.01 | 0.1520 | 0.3433 | 1026.24 |
| 5.0 | 2.47 | 2.02 | 0.0880 | 0.0948 | 1221.71 |
| 20.0 | 2.73 | 2.01 | 0.0280 |  | 1196.66 |
| 50.0 | 3.76 | 2.00 | 0.0120 |  | 1091.10 |

\begin{itemize}
\item ~\(\mu\)  = -0.86\(\mu \geq 5\)
\item rise/desc ~ \(\mu\)  = -1.05
\item 2.0Van der Pol
\item \textbf{Poincar } 6.3 (\(\mu=0.1\))  34.7 (\(\mu=20\))
\end{itemize}


\#\#

\subsubsection{D52}


1. \textbf{-}
\begin{itemize}
\item \(\mu\) \(\mu\)
\item -0.86  -1
\end{itemize}


2. ****
\begin{itemize}
\item Van der Pol  318.6
\item 0
\end{itemize}

-

3. \textbf{- trade-off}
\begin{itemize}
\item alpha
\item alpha  1.5
\end{itemize}


4. \textbf{alpha-mu }
\begin{itemize}
\item alpha-mu
\item pre-planning
\end{itemize}


\#\#

Van der Pol  \(\mu\) \textbf{-}
\begin{itemize}
\item \(\mu = 0\) \(f = 1/(2\pi) \approx 0.16\)
\item \(\mu > 0\) \(\mu\)  2.0
\item \(\mu \gg 1\)relaxation oscillation\(f \propto 1/\mu\)
\end{itemize}


\(\alpha\)

\#\#

\begin{itemize}
\item v3: `chimera_lab/d52-frequency-dynamics-redesign-v3.py`
\item v5: `chimera_lab/d52-frequency-dynamics-redesign-v5-vanderpol.py`
\item v5 results: `chimera_lab/d52-redesign-v5-vanderpol-results.json`
\item v5-F: `chimera_lab/d52-f-phase-diagram-relaxation.py`
\item v5-F results: `chimera_lab/d52-f-phase-diagram-results.json`
\end{itemize}


\begin{thebibliography}{99}
\bibitem{ref1} Shannon, C.E. (1948). A Mathematical Theory of Communication.
\bibitem{ref2} Cover, T.M. (2006). Elements of Information Theory.
\end{thebibliography}

\end{document}